\documentclass[10pt,a4paper]{amsart}
\usepackage{amsmath,amssymb,amsthm,graphicx}
% Retained minimal preamble for the Stacks Project chapter more-algebra.tex (More on Algebra).
% The chapter file pinned-original-source.tex carries only \input{preamble} on its first line, so no
% preamble text lives inside the chapter itself.  The Stacks Project uses one single shared preamble for
% every chapter; the snapshot stages it identically as sources/stacks-chow/preamble.tex and
% sources/stacks-etale-cohomology/preamble.tex.  The block below is copied mechanically, line for line,
% from sources/stacks-chow/preamble.tex lines 14-18 (the Xy-pic setup) and 154-217 (the native theorem
% declarations together with the mathematical macros).  Nothing was added, removed or redefined; the
% wrapper class and the amsmath/amssymb/amsthm/graphicx base packages are supplied by the assembler.
% Chapter: more-algebra.tex -- More on Algebra.
% For commutative diagrams we use Xy-pic
\usepackage[all]{xy}

% We use 2cell for 2-commutative diagrams.
\xyoption{2cell}

% Theorem environments.
%
\theoremstyle{plain}
\newtheorem{theorem}[subsection]{Theorem}
\newtheorem{proposition}[subsection]{Proposition}
\newtheorem{lemma}[subsection]{Lemma}

\theoremstyle{definition}
\newtheorem{definition}[subsection]{Definition}
\newtheorem{example}[subsection]{Example}
\newtheorem{exercise}[subsection]{Exercise}
\newtheorem{situation}[subsection]{Situation}

\theoremstyle{remark}
\newtheorem{remark}[subsection]{Remark}
\newtheorem{remarks}[subsection]{Remarks}

\numberwithin{equation}{subsection}

% Macros
%
\def\lim{\mathop{\mathrm{lim}}\nolimits}
\def\colim{\mathop{\mathrm{colim}}\nolimits}
\def\Spec{\mathop{\mathrm{Spec}}}
\def\Hom{\mathop{\mathrm{Hom}}\nolimits}
\def\Ext{\mathop{\mathrm{Ext}}\nolimits}
\def\SheafHom{\mathop{\mathcal{H}\!\mathit{om}}\nolimits}
\def\SheafExt{\mathop{\mathcal{E}\!\mathit{xt}}\nolimits}
\def\Sch{\mathit{Sch}}
\def\Mor{\mathop{\mathrm{Mor}}\nolimits}
\def\Ob{\mathop{\mathrm{Ob}}\nolimits}
\def\Sh{\mathop{\mathit{Sh}}\nolimits}
\def\NL{\mathop{N\!L}\nolimits}
\def\CH{\mathop{\mathrm{CH}}\nolimits}
\def\proetale{{pro\text{-}\acute{e}tale}}
\def\etale{{\acute{e}tale}}
\def\QCoh{\mathit{QCoh}}
\def\Ker{\mathop{\mathrm{Ker}}}
\def\Im{\mathop{\mathrm{Im}}}
\def\Coker{\mathop{\mathrm{Coker}}}
\def\Coim{\mathop{\mathrm{Coim}}}

% Boxtimes
%
\DeclareMathSymbol{\boxtimes}{\mathbin}{AMSa}{"02}

%
% Macros for moduli stacks/spaces
%
\def\QCohstack{\mathcal{QC}\!\mathit{oh}}
\def\Cohstack{\mathcal{C}\!\mathit{oh}}
\def\Spacesstack{\mathcal{S}\!\mathit{paces}}
\def\Quotfunctor{\mathrm{Quot}}
\def\Hilbfunctor{\mathrm{Hilb}}
\def\Curvesstack{\mathcal{C}\!\mathit{urves}}
\def\Polarizedstack{\mathcal{P}\!\mathit{olarized}}
\def\Complexesstack{\mathcal{C}\!\mathit{omplexes}}
% \Pic is the operator that assigns to X its picard group, usage \Pic(X)
% \Picardstack_{X/B} denotes the Picard stack of X over B
% \Picardfunctor_{X/B} denotes the Picard functor of X over B
\def\Pic{\mathop{\mathrm{Pic}}\nolimits}
\def\Picardstack{\mathcal{P}\!\mathit{ic}}
\def\Picardfunctor{\mathrm{Pic}}
\def\Deformationcategory{\mathcal{D}\!\mathit{ef}}

% --- appended 2026-10-12 (Chief): the project own definitions of the reference/slogan/history
% display environments, copied mechanically from the shared Stacks preamble lines 8-12
% (\usepackage{verbatim} plus the three \newenvironment lines, which the chapter preamble keeps
% verbatim from the project preamble). Nothing added, removed or redefined: these environments
% carry no mathematical content and the project itself typesets them as comments. Without these
% three lines a counted statement whose own text uses them fails with Environment slogan undefined.
\usepackage{verbatim}
\newenvironment{reference}{\comment}{\endcomment}
\newenvironment{slogan}{\comment}{\endcomment}
\newenvironment{history}{\comment}{\endcomment}
\usepackage{cleveref}
\crefname{originallemma}{original lemma}{original lemmas}
\Crefname{originallemma}{Original lemma}{Original lemmas}
\makeatletter
\newlabel{derived-lemma-subcategory-right-acyclics}{{05T8}{source}{}{sourcepointer-1}{}}
\newlabel{derived-lemma-subcategory-right-acyclics@cref}{{[originallemma][0][]05T8}{[1][0][]source}{}{}{}}
\makeatother
\crefname{originallemma}{original lemma}{original lemmas}
\Crefname{originallemma}{Original lemma}{Original lemmas}
\makeatletter
\newlabel{derived-lemma-unbounded-right-derived}{{07K7}{source}{}{sourcepointer-2}{}}
\newlabel{derived-lemma-unbounded-right-derived@cref}{{[originallemma][0][]07K7}{[1][0][]source}{}{}{}}
\makeatother
\crefname{originallemma}{original lemma}{original lemmas}
\Crefname{originallemma}{Original lemma}{Original lemmas}
\makeatletter
\newlabel{derived-proposition-enough-acyclics}{{05TA}{source}{}{sourcepointer-3}{}}
\newlabel{derived-proposition-enough-acyclics@cref}{{[originallemma][0][]05TA}{[1][0][]source}{}{}{}}
\makeatother
\crefname{originallemma}{original lemma}{original lemmas}
\Crefname{originallemma}{Original lemma}{Original lemmas}
\makeatletter
\newlabel{homology-lemma-Mittag-Leffler}{{02N1}{source}{}{sourcepointer-4}{}}
\newlabel{homology-lemma-Mittag-Leffler@cref}{{[originallemma][0][]02N1}{[1][0][]source}{}{}{}}
\makeatother
\crefname{originallemma}{original lemma}{original lemmas}
\Crefname{originallemma}{Original lemma}{Original lemmas}
\makeatletter
\newlabel{homology-section-inverse-systems}{{02MY}{source}{}{sourcepointer-5}{}}
\newlabel{homology-section-inverse-systems@cref}{{[originallemma][0][]02MY}{[1][0][]source}{}{}{}}
\makeatother
\title{The right derived functor of the inverse limit on N-indexed abelian groups has the standard two-term representation and vanishes above degree one}
\author{The Stacks Project Authors}
\date{stacks-project:more-algebra.tex:local-snapshot-20261012; GNU FDL 1.2 or later}
\begin{document}
\maketitle
\setcounter{subsection}{1}
\noindent\textit{Editorial note.} All delivered mathematical text is the author's own, copied from the pinned Stacks Project chapter file more-algebra.tex; no statement or proof step was rewritten, reordered or removed.\par\medskip
\noindent\textit{Editorial note.} The chapter file pinned-original-source.tex carries only \textbackslash{}input\{preamble\} on its first line, so no preamble text lives inside the chapter. The retained file sources/stacks-more-algebra/retained-preamble.tex is a mechanical, line-for-line copy of the project's single shared preamble as staged in the snapshot; nothing was added, removed or redefined and no package was added.\par\medskip
\noindent\textit{Editorial note.} The source URL is the source-level Stacks Project URL, as required when more than one card is produced from one source; the per-claim tag is carried by claim\_key.\par\medskip
\noindent\textit{Editorial note.} The counted statement span includes the author's numbered equation labelling the derived functor, delivered verbatim. The author writes 'of R\textasciicircum{}1 lim\_n K\textasciicircum{}p\_n = 0'; and the six references inside the counted proof span name the standard structural facts whose applications stay inside it.\par\medskip
\noindent\textit{Allowed standard prerequisites (editorial).}\begin{itemize}
\item The right-acyclic-subcategory criterion for the existence of a right derived functor. invoked at proof lines 22982-22986. Lets the explicitly built acyclics establish the derived functor and its degree bounds.
\item Mittag-Leffler inverse systems and the vanishing of their first derived limit. invoked at proof lines 22983-22984 and 22992. Is the acyclicity statement that the constructed subobject argument delivers.
\item The unbounded right derived functor lemma and its proof of the remaining assertions. invoked at proof lines 22994-22999. Extends the derived functor to the unbounded derived category and supplies properties (4) and (5).
\end{itemize}\medskip
\ifdefined\hypertarget\hypertarget{sourcepointer-1}{}\fi
\noindent\textit{Source reference.} Original lemma 05T8: focused-seeds/AGP/sources/stacks-snapshot-20261011/tags.txt, source lines 7472--7472. Cross-chapter navigation pointer into the Stacks Project; the standard fact it supplies is applied inside the counted proof span and the argument resting on it stays there.\par\medskip
\ifdefined\hypertarget\hypertarget{sourcepointer-2}{}\fi
\noindent\textit{Source reference.} Original lemma 07K7: focused-seeds/AGP/sources/stacks-snapshot-20261011/tags.txt, source lines 9647--9647. Cross-chapter navigation pointer into the Stacks Project; the standard fact it supplies is applied inside the counted proof span and the argument resting on it stays there.\par\medskip
\ifdefined\hypertarget\hypertarget{sourcepointer-3}{}\fi
\noindent\textit{Source reference.} Original lemma 05TA: focused-seeds/AGP/sources/stacks-snapshot-20261011/tags.txt, source lines 7474--7474. Cross-chapter navigation pointer into the Stacks Project; the standard fact it supplies is applied inside the counted proof span and the argument resting on it stays there.\par\medskip
\ifdefined\hypertarget\hypertarget{sourcepointer-4}{}\fi
\noindent\textit{Source reference.} Original lemma 02N1: focused-seeds/AGP/sources/stacks-snapshot-20261011/tags.txt, source lines 3464--3464. Cross-chapter navigation pointer into the Stacks Project; the standard fact it supplies is applied inside the counted proof span and the argument resting on it stays there.\par\medskip
\ifdefined\hypertarget\hypertarget{sourcepointer-5}{}\fi
\noindent\textit{Source reference.} Original lemma 02MY: focused-seeds/AGP/sources/stacks-snapshot-20261011/tags.txt, source lines 3461--3461. Cross-chapter navigation pointer into the Stacks Project; the standard fact it supplies is applied inside the counted proof span and the argument resting on it stays there.\par\medskip
\begin{lemma}
\label{lemma-compute-Rlim}
The functor $\lim : \textit{Ab}(\mathbf{N}) \to \textit{Ab}$
has a right derived functor
\begin{equation}
\label{equation-Rlim}
R\lim : D(\textit{Ab}(\mathbf{N})) \longrightarrow D(\textit{Ab})
\end{equation}
As usual we set $R^p\lim(K) = H^p(R\lim(K))$. Moreover, we have
\begin{enumerate}
\item for any $(A_n)$ in $\textit{Ab}(\mathbf{N})$ we have
$R^p\lim A_n = 0$ for $p > 1$,
\item the object $R\lim A_n$ of $D(\textit{Ab})$ is represented
by the complex
$$
\prod A_n \to \prod A_n,\quad (x_n) \mapsto (x_n - f_{n + 1}(x_{n + 1}))
$$
sitting in degrees $0$ and $1$,
\item if $(A_n)$ is ML, then $R^1\lim A_n = 0$, i.e., $(A_n)$
is right acyclic for $\lim$,
\item every $K^\bullet \in D(\textit{Ab}(\mathbf{N}))$ is quasi-isomorphic
to a complex whose terms are right acyclic for $\lim$, and
\item if each $K^p = (K^p_n)$ is right acyclic for $\lim$, i.e.,
of $R^1\lim_n K^p_n = 0$, then $R\lim K$ is represented by the
complex whose term in degree $p$ is $\lim_n K_n^p$.
\end{enumerate}
\end{lemma}

\begin{proof}
Let $(A_n)$ be an arbitrary inverse system. Let $(B_n)$ be the inverse
system with
$$
B_n = A_n \oplus A_{n - 1} \oplus \ldots \oplus A_1
$$
and transition maps given by projections. Let $A_n \to B_n$ be given
by $(1, f_n, f_{n - 1} \circ f_n, \ldots, f_2 \circ \ldots \circ f_n)$
where $f_i : A_i \to A_{i - 1}$ are the transition maps.
In this way we see that every inverse system is a subobject of a
ML system (Homology, Section \ref{homology-section-inverse-systems}).
It follows from
Derived Categories, Lemma \ref{derived-lemma-subcategory-right-acyclics}
using Homology, Lemma \ref{homology-lemma-Mittag-Leffler}
that every ML system is right acyclic for $\lim$, i.e., (3) holds.
This already implies that $RF$ is defined on $D^+(\textit{Ab}(\mathbf{N}))$,
see Derived Categories, Proposition \ref{derived-proposition-enough-acyclics}.
Set $C_n = A_{n - 1} \oplus \ldots \oplus A_1$ for $n > 1$ and
$C_1 = 0$ with transition maps given by projections as well.
Then there is a short exact sequence of inverse systems
$0 \to (A_n) \to (B_n) \to (C_n) \to 0$ where $B_n \to C_n$
is given by $(x_i) \mapsto (x_i - f_{i + 1}(x_{i + 1}))$.
Since $(C_n)$ is ML as well, we conclude that (2) holds
(by proposition reference above) which also implies (1).
Finally, this implies by Derived Categories, Lemma
\ref{derived-lemma-unbounded-right-derived}
that $R\lim$ is in fact defined on all of $D(\textit{Ab}(\mathbf{N}))$.
In fact, the proof of Derived Categories, Lemma
\ref{derived-lemma-unbounded-right-derived}
proceeds by proving assertions (4) and (5).
\end{proof}

\end{document}
